% ============================================================================
% Section « Pistes explorées et difficultés rencontrées » — à insérer entre
% la section 4 (Justification des choix méthodologiques) et la section 5
% (Valeur créée et limites). La numérotation LaTeX se recale toute seule.
% Aucun package particulier requis.
% ============================================================================

\section{Pistes explorées et difficultés rencontrées}

La solution présentée dans les sections précédentes est le produit d'une
démarche itérative : plusieurs pistes ont été ouvertes, certaines abandonnées,
et deux erreurs méthodologiques significatives ont été commises puis corrigées.
Nous les documentons ici, car elles ont directement façonné les garde-fous du
livrable final.

\subsection{Pistes ouvertes puis écartées}

Trois pistes envisagées lors de l'analyse préliminaire n'ont pas été retenues.
La première était un modèle de \emph{durée de blocage} du trafic, exploitant
l'écart entre l'heure de début et de fin de chaque accident ; elle a été
abandonnée lorsque l'exploration a montré que 44~\% des coordonnées de fin
étaient manquantes, jetant un doute sérieux sur la fiabilité des horodatages
de fin eux-mêmes. La deuxième concernait le périmètre géographique : nous
avons hésité entre restreindre l'étude au comté de Los Angeles (lisibilité
maximale), à un État médian comme la Caroline du Sud (volume modeste mais
équilibré), ou à la Californie entière. Le débat a été tranché par un
argument de généricité : entraîner sur l'ensemble des États-Unis et
paramétrer la démonstration par une simple constante de territoire, la
Californie ne servant que de terrain d'illustration. La troisième piste était
une chaîne de modélisation entièrement portée par Spark~MLlib ; elle a été
écartée au profit de scikit-learn sur un échantillon plafonné et documenté,
seul moyen de mener dans le temps imparti une comparaison riche de cinq
modèles avec les métriques et les outils d'explicabilité souhaités. MLlib
reste mentionné comme voie de passage à l'échelle.

\subsection{Une impasse instructive : le modèle tautologique}

La première tentative de réponse au « Quoi » constitue l'erreur la plus
instructive du projet. Nous avions entraîné un modèle de régression
(forêt aléatoire) prédisant la \emph{charge} d'une zone — la somme des
sévérités de ses accidents — à partir des comptages d'accidents impliquant
chaque équipement dans cette même zone. Le coefficient de détermination était
excellent, et pour cause : la cible était une fonction quasi arithmétique des
variables d'entrée, puisque les comptages par équipement \emph{composent} la
charge. Le modèle n'apprenait pas un phénomène routier ; il reconstituait une
identité comptable, et ses importances de variables ne décrivaient que la
composition du territoire.

Cette impasse a produit deux garde-fous permanents. D'une part, le changement
d'unité d'analyse : le modèle final travaille au niveau de l'accident
individuel, dont la gravité n'est pas une fonction de ses caractéristiques
d'infrastructure. D'autre part, des tests automatisés anti-fuite, exécutés à
chaque exécution du notebook, qui vérifient qu'aucune variable d'entrée ne
contient la cible sous quelque forme que ce soit et que le découpage temporel
train/test est étanche.

\subsection{Le piège des effets bruts, en deux temps}

La seconde correction majeure concerne la chaîne d'explication elle-même, et
s'est jouée en deux temps. Le premier est décrit en
section~4.3 : les valeurs de Shapley globales suggéraient que feux et
passages piétons aggravaient la gravité, artefact du biais de sélection des
équipements. Le second est plus subtil : notre première version de la
prescription par zone reposait sur des écarts de sévérité moyens calculés
directement dans les données — alors même que l'argument central du recours
au modèle était sa capacité à contrôler les facteurs de confusion. La
prescription contredisait donc silencieusement sa propre justification. Nous
avons résolu cette incohérence en faisant piloter la recommandation par
l'effet ajusté conditionnel issu du calcul-G, l'écart brut observé n'étant
conservé que comme contre-vérification. C'est également ce qui a motivé
l'ajout des rapports de cotes de la régression logistique, avec leurs
intervalles de confiance : deux lectures indépendantes qui concordent sur
dix équipements sur douze constituent une garantie qu'aucune des deux ne
fournit isolément.

\subsection{Difficultés d'ingénierie}

Deux difficultés techniques méritent mention. La première est la volumétrie :
avec près de huit millions de lignes, chaque erreur de conception de la
pipeline Spark se paie en dizaines de minutes de recalcul ; la discipline de
mise en cache des jeux intermédiaires et de plafonnement explicite des
échantillons a été acquise par la pratique. La seconde, plus sournoise, est
la dépendance au service Overpass d'OpenStreetMap pour l'exposition
routière : requêtes bloquantes sans délai de garde, limitations de débit
agressives et imprévisibles, et un piège de configuration — la clé de cache
du client dépend du délai d'expiration configuré, si bien qu'un cache
constitué avec un réglage différent devient invisible. La version finale
encapsule cette étape dans un budget de temps global avec dégradation
contrôlée : en l'absence de réponse du service, le classement par volume est
conservé et la limite documentée, de sorte que l'exécution complète du
notebook aboutit dans tous les cas. Le cache local est en outre livré avec le
dépôt, ce qui permet de rejouer l'exposition hors-ligne sur quarante-deux des
cinquante zones.

Enfin, la suite de tests a démontré son utilité dès sa première exécution en
détectant une régression de contenu : une édition antérieure du notebook
avait amputé, sans que personne ne s'en aperçoive, la section d'analyse
critique. L'anecdote illustre concrètement pourquoi un livrable, même sous
forme de notebook, gagne à être protégé par un protocole de validation
automatisé.

% ============================================================================
% CORRECTIONS À APPLIQUER AILLEURS DANS LE RAPPORT (sur-affirmations vs code)
% ============================================================================
% 1) §3.1 — REMPLACER « les valeurs manquantes des variables météorologiques
%    (température, humidité, visibilité) ont été imputées par localité »
%    PAR : « les valeurs manquantes des variables météorologiques ont été
%    imputées par la médiane, choix robuste aux valeurs extrêmes »
% 2) §3.1 — REMPLACER « les enregistrements aux coordonnées géographiques
%    aberrantes, manquantes ou hors des frontières attendues ont été
%    systématiquement exclus »
%    PAR : « les enregistrements sans coordonnées ou sans gravité renseignée
%    ont été systématiquement exclus »
% 3) §3.3 — REMPLACER « L'indice de risque retenu est le nombre d'accidents
%    rapporté au kilométrage de voirie, pondéré par le type de route »
%    PAR : « Les zones sont d'abord classées par leur charge (somme des
%    sévérités) ; pour les zones de tête, l'intensité par kilomètre de voirie
%    (via OpenStreetMap) est ensuite calculée afin de distinguer les zones
%    réellement dangereuses des zones simplement très pourvues en voirie »
% 4) §4.3 — REMPLACER « l'impact de chaque aménagement (via le calcul-G) est
%    mesuré uniquement au sein de sous-ensembles de zones présentant un
%    niveau comparable d'aléa local et de complexité. Cette stratification… »
%    PAR : « l'impact de chaque aménagement (via le calcul-G) est mesuré
%    uniquement sur les accidents survenus en présence de l'aléa local
%    dominant de la zone, le contexte (météo, heure, État) étant maintenu
%    constant par le modèle. Ce conditionnement… »
% 5) §1 — « près de sept millions d'enregistrements » → « près de huit
%    millions d'enregistrements (7,73 millions) »
% 6) Figure 2 — la légende annonce une comparaison avant/après normalisation
%    OSM mais la figure montre la distribution des scores : remplacer par la
%    figure de concentration du notebook final (courbe cumulée + volume ×
%    gravité) ou corriger la légende.
% ============================================================================
